\chapter{Manuel --- Secr\'etariat P\'edagogique}
\label{chap:secretariat}

\begin{rolebox}{secretariat-pedagogique}
  Ce chapitre est destin\'e aux agents du secr\'etariat p\'edagogique. Ce profil est orient\'e \textbf{consultation et mise \`a jour des contacts fonctionnels}. Il ne permet pas de g\'erer des utilisateurs, de cr\'eer des affectations ou de g\'en\'erer des organigrammes.
\end{rolebox}

\section{Vue d'ensemble du profil}

Le profil \texttt{secretariat-pedagogique} dispose des fonctionnalit\'es suivantes\,:

\begin{itemize}
  \item \textbf{Consultation de l'annuaire} : recherche de responsables, formations, structures et secr\'etariats.
  \item \textbf{Consultation des fiches structures} : acc\`es aux d\'etails complets des structures.
  \item \textbf{Mise \`a jour des contacts de r\^ole} : modification des informations de contact fonctionnel li\'ees aux affectations (courriel fonctionnel, t\'el\'ephone, bureau).
  \item \textbf{Consultation des organigrammes} g\'en\'er\'es (lecture seule).
  \item \textbf{Signalements} : signalement d'erreurs sur les donn\'ees.
  \item \textbf{Modification du profil personnel}.
\end{itemize}

\begin{notebox}[Ce qui n'est PAS disponible]
  Le secr\'etariat p\'edagogique ne peut \textbf{pas}\,:
  \begin{itemize}
    \item Cr\'eer, modifier ou supprimer des comptes utilisateurs.
    \item Cr\'eer ou supprimer des affectations.
    \item G\'en\'erer de nouveaux organigrammes.
    \item Importer ou exporter des donn\'ees.
    \item G\'erer des d\'el\'egations.
    \item G\'erer les ann\'ees acad\'emiques.
  \end{itemize}
\end{notebox}

\section{Tableau de bord}

\screenshot{img_dashboard_secretariat}{Tableau de bord -- Secr\'etariat P\'edagogique}{Capturer le tableau de bord apr\`es connexion avec un compte secretariat-pedagogique. Montrer les tuiles de raccourcis limit\'ees \`a ce profil (Rechercher, Organigramme, Signalements, Profil) et l'ann\'ee courante.}

\section{Consultation de l'annuaire}

L'annuaire est la fonctionnalit\'e principale pour le secr\'etariat. Il permet de retrouver rapidement toute personne, formation ou structure de l'\'etablissement.

\begin{steps}
  \item Cliquer sur \textbf{Rechercher} dans le menu.
  \item Choisir l'onglet correspondant \`a votre recherche.
  \item Saisir un texte libre ou utiliser les filtres.
\end{steps}

\subsection{Onglet Responsables}

Permet de trouver une personne par\,: nom, pr\'enom, courriel, r\^ole, composante, d\'epartement, mention, parcours ou niveau.

\screenshot{img_annuaire_responsables_secretariat}{Annuaire -- Onglet Responsables (Secr\'etariat)}{Capturer la page Recherche avec l'onglet Responsables actif. Montrer la barre de recherche texte, les filtres hi\'erarchiques (composante, d\'epartement, mention, parcours, niveau) et les filtres de r\^ole, et un exemple de r\'esultats avec des cartes ou lignes de r\'esultat.}

\subsection{Onglet Formations}

Permet de rechercher une formation par son intitul\'e, son type de dipl\^ome, sa composante ou son d\'epartement.

\screenshot{img_annuaire_formations_secretariat}{Annuaire -- Onglet Formations (Secr\'etariat)}{Capturer l'onglet Formations de la recherche avec des r\'esultats affich\'es. Montrer les filtres disponibles (type de dipl\^ome, composante, d\'epartement) et les cartes de r\'esultat avec intitul\'e de formation et responsables.}

\subsection{Onglet Structures}

Permet de retrouver une structure (composante, d\'epartement, mention, parcours, service) par son nom ou sa hi\'erarchie.

\screenshot{img_annuaire_structures_secretariat}{Annuaire -- Onglet Structures (Secr\'etariat)}{Capturer l'onglet Structures de l'annuaire avec des r\'esultats. Montrer les filtres et le tableau de r\'esultats (nom de structure, type, parent).}

\subsection{Onglet Secr\'etariats}

Permet de retrouver les contacts de secr\'etariat d'une composante ou d'un d\'epartement.

\screenshot{img_annuaire_secretariats_secretariat}{Annuaire -- Onglet Secr\'etariats (Secr\'etariat)}{Capturer l'onglet Secr\'etariats avec des r\'esultats. Montrer les informations affich\'ees\,: nom de la personne, structure rattach\'ee, courriel fonctionnel, t\'el\'ephone, bureau.}

\section{Consultation des fiches structures}

Le secr\'etariat peut consulter les fiches d\'etaill\'ees de toutes les structures.

\begin{steps}
  \item Rechercher la structure dans l'annuaire (onglet \textbf{Structures}).
  \item Cliquer sur la structure pour ouvrir sa fiche.
  \item La fiche affiche\,: identification, hi\'erarchie, coordonn\'ees, responsables, secr\'etariat, sous-structures.
\end{steps}

\screenshot{img_fiche_structure_secretariat}{Fiche structure -- Vue Secr\'etariat}{Capturer la fiche d\'etaill\'ee d'une composante ou d'un d\'epartement. Toutes les sections doivent \^etre visibles\,: titre, type, ann\'ee, parent, coordonn\'ees, liste des responsables (nom, r\^ole, courriel), liste du secr\'etariat, sous-structures directes.}

\section{Mise \`a jour des contacts de r\^ole}

Le secr\'etariat p\'edagogique peut mettre \`a jour les informations de contact fonctionnel associ\'ees aux affectations (courriel fonctionnel, t\'el\'ephone de service, bureau).

\begin{warnbox}[Distinction contact personnel / contact fonctionnel]
  Le \textbf{contact personnel} (courriel priv\'e, t\'el\'ephone personnel) est g\'er\'e via \og Mon Profil \fg{}. \\
  Le \textbf{contact de r\^ole} (courriel fonctionnel de service, t\'el\'ephone du secr\'etariat, bureau du service) est g\'er\'e via la rubrique \textbf{Responsables} \`a travers les affectations.
\end{warnbox}

\subsection{Trouver l'affectation \`a mettre \`a jour}

\begin{steps}
  \item Naviguer vers \textbf{Responsables}.
  \item Rechercher la personne ou filtrer par structure.
  \item Localiser l'affectation dont les coordonn\'ees doivent \^etre mises \`a jour.
\end{steps}

\subsection{Modifier le contact de r\^ole}

\begin{steps}
  \item Dans la ligne ou fiche de l'affectation, cliquer sur \textbf{Contact de r\^ole} (ou ic\^one crayon).
  \item Modifier les champs souhait\'es\,:
    \begin{itemize}
      \item \textit{Courriel fonctionnel} : courriel de service (ex. : \texttt{secretariat.droit@univ.fr}).
      \item \textit{T\'el\'ephone de service} : num\'ero du bureau ou du secr\'etariat.
      \item \textit{Bureau / Localisation} : num\'ero de bureau, b\^atiment, campus.
    \end{itemize}
  \item Cliquer sur \textbf{Enregistrer}.
\end{steps}

\screenshot{img_contact_role_secretariat}{Modification du contact de r\^ole -- Secr\'etariat}{Capturer le formulaire d'\'edition du contact de r\^ole tel qu'il appara\^it pour un secr\'etariat. Montrer les trois champs (courriel fonctionnel, t\'el\'ephone, bureau) avec leurs labels et un exemple de valeur saisie.}

\begin{tipbox}[Cas d'usage typique]
  En d\'ebut d'ann\'ee acad\'emique ou lors d'un changement de locaux, le secr\'etariat met \`a jour les coordonn\'ees fonctionnelles de toutes les affectations de son service afin que l'annuaire reste \`a jour pour l'ensemble des utilisateurs.
\end{tipbox}

\section{Consultation des organigrammes}

Le secr\'etariat peut consulter les organigrammes d\'ej\`a g\'en\'er\'es mais ne peut pas en cr\'eer de nouveaux.

\begin{steps}
  \item Naviguer vers \textbf{Organigramme}.
  \item Acc\'eder \`a la section \og Biblioth\`eque \fg{}.
  \item Filtrer par structure ou type de vue.
  \item Cliquer sur \textbf{Voir en structures} ou \textbf{Voir en personnes}.
\end{steps}

\screenshot{img_organigramme_consultation_secretariat}{Consultation d'organigramme -- Secr\'etariat}{Capturer l'organigramme en consultation (vue personnes ou vue structures). Montrer que les boutons de g\'en\'eration et de gel ne sont pas visibles (ou sont d\'esactiv\'es) pour ce profil.}

\begin{notebox}
  Le secr\'etariat peut consulter les organigrammes de toutes les structures, pas uniquement ceux de sa propre composante.
\end{notebox}

\section{Signalements}

Le secr\'etariat peut signaler des erreurs sur les donn\'ees.

\begin{steps}
  \item Naviguer vers \textbf{Signalements}.
  \item Cliquer sur \textbf{Cr\'eer un signalement}.
  \item S\'electionner le type d'erreur\,:
    \begin{itemize}
      \item \textbf{Nom} : erreur sur le nom d'une personne.
      \item \textbf{Pr\'enom} : erreur sur le pr\'enom.
      \item \textbf{T\'el\'ephone} : num\'ero incorrect ou manquant.
      \item \textbf{Bureau} : localisation incorrecte.
      \item \textbf{Courriel} : adresse courriel incorrecte.
      \item \textbf{R\^ole} : affectation incorrecte ou manquante.
      \item \textbf{Autre} : tout autre type d'erreur.
    \end{itemize}
  \item Identifier la \textit{personne} ou la \textit{structure} concern\'ee.
  \item D\'ecrire le probl\`eme dans le champ \textit{Description}.
  \item Cliquer sur \textbf{Envoyer}.
\end{steps}

\screenshot{img_signalements_secretariat}{Signalements -- Secr\'etariat P\'edagogique}{Capturer la page Signalements avec\,: (1) le formulaire de cr\'eation (menu d\'eroulant du type d'erreur, champ de recherche de la personne concern\'ee, champ description), et (2) la liste des signalements d\'ej\`a soumis avec leurs statuts.}

\begin{tipbox}[Suivi des signalements]
  Apr\`es soumission, vous pouvez suivre le statut de vos signalements depuis la page \textbf{Signalements}. Une notification vous est envoy\'ee lorsque votre signalement est trait\'e ou cl\^otur\'e.
\end{tipbox}

\section{Mon Profil}

Le secr\'etariat peut consulter et mettre \`a jour ses propres coordonn\'ees.

\begin{steps}
  \item Cliquer sur \textbf{Mon Profil} dans le menu.
  \item Cliquer sur \textbf{Modifier} dans la section coordonn\'ees.
  \item Mettre \`a jour\,: \textit{Courriel secondaire}, \textit{T\'el\'ephone}, \textit{Bureau}.
  \item Cliquer sur \textbf{Enregistrer}.
\end{steps}
